\Needspace{9\baselineskip}
\subsection{资产配置收益归因与下跌日表现}
\label{sec:brinson}
\subsubsection{仓位、风险贡献与收益归因的分析对象}
仓位管理决定资本投向，风险贡献衡量各项持仓对组合风险的影响，收益归因则解释组合相对基准的实际收益差。本文以权重和协方差计算资产层面的波动风险贡献，并以Brinson--Fachler框架核算配置收益。Barra采用因子暴露、因子协方差和特异风险刻画组合风险，其因子收益归因还需要相应的因子收益序列\cite{msci2026factor}。因此，本文以下归因的分析单位是ETF资产组，结论对应可观察的持仓决策。

Brinson与Fachler提出了按类别分解配置和选择作用的框架\cite{brinson1985measuring}；Brinson、Hood与Beebower进一步讨论了投资政策、配置和选择决策的绩效评价\cite{brinson1986determinants}。本文采用Brinson--Fachler的类别配置口径，并依据Bacon的专业综述分别报告配置、选择与交互项\cite{bacon2019attribution}。这一框架适合回答：相对同一比较基准，多配或少配某类资产，贡献了多少实际收益差。

\subsubsection{计算口径及现金配置的含义}
设$w^P_{g,t}$与$w^B_{g,t}$分别为组合和基准在当日收益发生前的第$g$组权重，$r^P_{g,t}$与$r^B_{g,t}$为组内按资本加权的当日收益，$R^B_t=\sum_g w^B_{g,t}r^B_{g,t}$为基准毛收益。配置、组内选择与交互效应分别为
\begin{align}
A_{g,t}&=(w^P_{g,t}-w^B_{g,t})(r^B_{g,t}-R^B_t),\\
S_{g,t}&=w^B_{g,t}(r^P_{g,t}-r^B_{g,t}),\\
I_{g,t}&=(w^P_{g,t}-w^B_{g,t})(r^P_{g,t}-r^B_{g,t}).
\end{align}
满额投资使两组权重各自加总为一，故$\sum_g(A_{g,t}+S_{g,t}+I_{g,t})=R^P_t-R^B_t$。将两条路径的费用率分别记作$c^P_t$、$c^B_t$，另列费用效应$C_t=c^B_t-c^P_t$后，各项之和等于当日净收益差。

配置项衡量类别权重偏离带来的相对贡献。若组合超配货币ETF，而且货币组收益高于当日基准整体收益，则货币配置项为正。这并不要求货币ETF自身取得很高收益。例如，在一个仅含现金与股票的假设单期中，现金收益为零、股票下跌10\%，基准持有10\%现金、组合持有50\%现金，两者收益分别为$-9\%$与$-5\%$。组合领先4个百分点，其中现金配置项为3.6个百分点，股票低配项为0.4个百分点。这一例子区分了现金自身的收益贡献、现金超配的归因贡献和整个仓位调整的净效果。

组内选择项反映同一类别中ETF权重构成的差异；本文所称“选择”对应基金配置，不涉及穿透基金后的个股选择。交互项记录类别权重偏离与组内收益偏离同时出现的影响。本文将其独立列示，使三项之和可以逐日复核。该分解属于给定基准下的会计恒等式，关于配置规则作用的实验依据由前述消融与风险匹配提供。

\subsubsection{资产分组、信息时序与核算结果}
为单独观察现金、利率和信用配置，本文在归因中设置货币、国债、信用债、可转债、境内权益、境外权益与商品相关七个互斥组。该分组用于解释收益，不改变资产池或优化规则。港股计入境外权益；商品相关组沿用项目工具映射，包含煤炭等与商品相关的权益工具，因此其解释保持在ETF类别层面。

基准分别采用Equal Weight及其事前风险匹配路径。计算使用与已公布日收益一致的权重：调仓日使用当日目标，其他日期使用收益漂移后的持仓。组内收益按当日组权重归一化。若组合某组权重为零，其组收益取原始等权基准的对应组收益；若风险匹配基准某组为空，同样采用原始等权组收益作为类别参照。这一约定保持空组解释一致并保留逐日恒等式。空组的参照收益以及资产分组会影响分项数值，因此本文固定上述口径，并同时核查合计收益差。全部持有资产在归因日期均有有效收益，重建毛收益与已公布路径一致。

表\ref{tab:brinson_summary}使用日均基点表示各项贡献，一个基点等于0.01个百分点。该统计是单日效应的算术平均，适合比较全样本和条件样本；复合累计收益及年化收益仍由前文净值路径计算。Bacon讨论了多期归因需要处理收益链接的问题\cite{bacon2019attribution}，本文据此将单日均值与复合绩效分别报告。

归因样本与主绩效中由相邻净值计算的收益样本对齐，2018年1月2日作为净值基期，归因从下一交易日开始。每日四项效应之和等于当日净收益差，因而各项均值之和也等于平均日净收益差。对这些均值乘以252只会得到算术年化差异，无法直接还原两项复合年化收益之差。表中各项独立四舍五入，显示值的加总可能有末位差异，逐日未舍入数值另行核验。



全样本共\attrDays{}个交易日，沪深300ETF当日收益小于零的观察有\attrDownDays{}个。下跌日按当期实现收益事后划分，用于解释已实现损益，与表\ref{tab:position_states}中依据历史信息划分的弱势状态口径不同。相对Equal Weight，主模型在下跌日的平均净收益差为\attrEqualNetDown{}个基点，全样本对应\attrEqualNetAll{}个基点；这一组合同时体现了下跌时的损失缓冲和较低风险配置的长期收益机会成本。相对风险匹配等权路径，全样本和下跌日的平均净收益差分别为\attrMatchedNetAll{}和\attrMatchedNetDown{}个基点。

\begin{table}[htbp]\centering
\caption{各资产组的配置效应：平均单日贡献（基点）}\label{tab:brinson_allocation}
\fitthesistable{\begin{tabular}{lrrrr}
\toprule
归因组 & 等权全样本 & 等权下跌日 & 匹配全样本 & 匹配下跌日\\\midrule
货币 & -0.34 & 7.04 & 0.38 & -7.06\\
国债 & -0.58 & 14.43 & 0.04 & 4.85\\
信用债 & -0.31 & 4.75 & -0.06 & 1.42\\
可转债 & -0.00 & 0.01 & 0.01 & -0.36\\
境内权益 & 0.17 & 13.35 & 0.03 & 1.70\\
境外权益 & -0.14 & -1.08 & -0.02 & -1.02\\
商品相关 & -0.21 & -1.94 & 0.33 & -0.66\\
\bottomrule
\end{tabular}}
\end{table}
\FloatBarrier

表\ref{tab:brinson_allocation}进一步显示，相对原始等权基准，下跌日的货币、国债及信用债配置均提供正的平均配置贡献，其中货币项为\attrEqualCashDown{}个基点。风险匹配等权路径本身配置了较多货币ETF，主模型相对该路径的货币项为\attrMatchedCashDown{}个基点，国债、信用债与境内权益配置项则为正。两组结果说明，现金和债券可以通过资本配置参与下行风险管理，而具体贡献的符号取决于基准已有的持仓。主模型相对风险匹配路径的收益来源还包括组内配置和交互效应，因而需要结合完整分解评价。

两个基准承担不同解释任务：原始等权显示风险规模和配置结构共同形成的收益差，风险匹配路径显示指定事前方差尺度下的差异。下跌日分组会突出低权益敞口的损失缓冲，表中结果应与全样本及实现风险同时阅读。沪深300ETF下跌日只是境内宽基下行的观察集合，不能代表所有市场的同步危机，也不构成策略提前预测下跌的证据。
\FloatBarrier
